\documentclass[11pt]{article}
\usepackage{amsmath,amssymb}
\newcommand{\Gen}[1]{\langle #1\rangle}
\title{Three finite certificates in the original Fi22 marking domain}
\author{}
\date{}
\begin{document}
\maketitle

We retain the accepted original degree $3510$ action of the inner simple
group $G=\mathrm{Fi}_{22}$, its procedural family of involutions, and the
original free-row identifications. In particular, membership in this same
$G$ and transitivity of its point action are borrowed model premises.
Membership merely in $\operatorname{Aut}(G)$ would not suffice.

\paragraph{Proposition.}
The original markings at free indices $310$ and $312$, with owners $652$
and $655$, generate $G$ and have the full rank-three intersection property.
The marking at free index $313$, with owner $656$, does not generate $G$.

\paragraph{Proof.}
For each marking write its actual ordered involutions as $(a,b,c)$ and put
\[
                 K=\Gen{a,b,c},\qquad R=\Gen{ab,bc}.
\]
The unchanged strict worker reconstructs the complete endpoint and middle
permutations. It checks nonidentity, involution and distinctness of the
three entries, as well as $ac=ca$. Its original-row checks identify the
entire $16$-byte witness row, not merely the free index or owner.

At each of indices $310$ and $312$, the separate checker verifies two
supplied signed words in $ab,bc$. One word has a closed cycle of exactly
$11$ distinct points, and the other has a closed cycle of exactly $13$
distinct points. Every closing edge is checked in the actual original
action. Hence $11$ and $13$ divide $|R|$.

The accepted maximal subgroup criterion says that no proper subgroup of
the named simple $\mathrm{Fi}_{22}$ has both prime divisors. It follows
that $R=G$, and therefore $K=G$. This criterion uses the published maximal
subgroup classification, with its correction; it is not a new recognition
theorem. The existing rank-three intersection theorem applies to the
checked structural hypotheses and $K=G$, giving the full intersection
property for both markings.

For index $313$, the supplied certificate is a set $S$ of $3159$ original
point labels. The strict checker verifies its distinct sorted entries and
all $2\cdot3159=6318$ positive rotation edges, establishing
\[
      \varnothing\ne S\subsetneq\Omega,\qquad S^{ab}\subseteq S,
      \qquad S^{bc}\subseteq S,
      \qquad |\Omega|=3510.
\]
Since these are permutations of a finite set, the containments are
equalities. Thus $R$ preserves the proper nonempty subset $S$ and is
intransitive on $\Omega$.

The subgroup $R$ is normal in $K$ and has index at most two: $K$ is
generated by $R$ and $b$, and conjugation by $b$ inverts both generators
$ab$ and $bc$. If $K=G$, simplicity and nonabelianness would force
$R=G$. That contradicts the borrowed transitivity of $G$. Hence $K<G$.
No failed search is used in this exclusion.
\hfill$\square$

\paragraph{Attribution and verification.}
The maximal subgroup classification is due to P. B. Kleidman and
R. A. Wilson, \emph{The maximal subgroups of $\mathrm{Fi}_{22}$},
Math. Proc. Cambridge Philos. Soc. \textbf{102} (1987), 17--23,
doi:10.1017/S0305004100067001, with the corrigendum in \textbf{103}
(1988), 383, doi:10.1017/S0305004100064963.
The retained argument for the two-prime criterion records the complete
classification cover and the correction's scope. The rank-three
intersection theorem and its proof are reused unchanged from
\texttt{rank\_three\_intersections.tex}.

The new work consists of these supplied finite certificates, their
separate strict checks and their applications to the stated original
markings. The model, domain identifications, classification and general
intersection theorem remain attributed earlier inputs. The exact
certificates, genuine producer and checker calls, Source, controls,
resource receipts and independent reviews are bound in the application
records. Ledger adoption is a separate step. These three entries do not
complete the manuscript, settle an open conjecture, constitute a Lean
formalization, or change a compression measurement.
\end{document}
